\documentclass[11pt]{article}
\usepackage[T1]{fontenc}
\usepackage{lmodern}
\usepackage[margin=1in]{geometry}
\usepackage{microtype}
\usepackage{amsmath}
\usepackage[backend=biber,style=authoryear]{biblatex}
\addbibresource{../../resources/references.bib}
\usepackage{hyperref}

\title{A Reproducible Writing Workflow}
\author{LaTeX for Everyone}
\date{\today}

\begin{document}
\maketitle
\begin{abstract}
This short teaching article demonstrates an abstract, numbered sections,
a cross-reference, and a bibliography. It proposes a writing workflow,
not an empirical research finding.
\end{abstract}

\section{Introduction}
\label{sec:intro}
Structured markup lets authors revise content without manually
renumbering headings \parencite{lamport1994}.

\section{Method}
Write the central argument first, organize supporting evidence, and
compile after small edits. Keep the source and build instructions
under version control.

\section{Discussion}
The approach in Section~\ref{sec:intro} still requires human review.
Successful compilation does not establish factual accuracy,
accessibility, or compliance with a journal's submission rules.

\section{Conclusion}
Use this document as a starting point, then replace its argument and
citations with your own original work.
\printbibliography
\end{document}
